\section{Methods and Correctness}
\label{sec:methods}

All conditions use one graph-search engine and differ only in when the nested
heuristic is evaluated.  Figure~\ref{fig:stages} shows the proposed staged ladder.

\begin{figure}[t]
  \centering
  \includegraphics[width=\columnwidth]{stage_schematic.pdf}
  \caption{The proposed three-level staged-evaluation ladder.  A non-goal
  \Staged{} state expands only after reaching \HFull{}.}
  \label{fig:stages}
\end{figure}

\paragraph{\Manhattan.}
The start and every accepted state receive only \HMan{} and expand on a valid
pop.  This is the cheap, weak baseline.

\paragraph{\EagerFull.}
Every accepted state receives \HFull{} immediately.  This avoids heuristic
promotion and requeue overhead but pays all 32 landmark terms even for states
that never expand.

\paragraph{\LazyFull.}
Following Lazy \Astar{}, the start is fully evaluated before insertion.
Other accepted states initially receive \HMan{}.  A valid non-goal pop at the
cheap stage computes \HFull{} and requeues the state with the same insertion
ordinal; a later full-stage pop expands it.  A promotion occurs even when the
numeric estimate does not increase.

\paragraph{\Staged.}
The start computes \HPrefix{} and then the remaining 28 terms before its
initial insertion.  Other states begin with \HMan{}.  Their first valid
non-goal pop computes \HPrefix{} and requeues; a second competitive pop
computes only the remaining terms and requeues; a full-stage pop expands.
This method cannot bypass the suffix if a state is to expand.

\subsection{Shared graph-search contract}

OPEN entries are ordered lexicographically by
\begin{equation}
 (f,-g,\textit{ordinal},\textit{row-major id},
  \textit{version},\textit{stage}).
 \label{eq:open}
\end{equation}
A successful relaxation receives a fresh ordinal; heuristic promotion
preserves it.  Per-state labels store the best $g$, parent, version, ordinal,
completed stage, and current $h$.  Stale entries are discarded.  Cached
heuristic stages survive a later strict better-$g$ relaxation because their
values do not depend on $g$.  CLOSED reopening is implemented even though
consistency made its count zero in this experiment.

On a valid pop, the goal test precedes expensive refinement.  Every heuristic
is zero at the goal, so refinement could not raise its key.  A non-goal state
enumerates successors only at its condition's final stage.  We define
\emph{expanded} accordingly; the returned goal is not counted as expanded and
does not enter the expansion digest.

\subsection{Optimality and trace equivalence}

At every stage, $f=g+h$ is a lower bound on the cost of a solution through the
entry.  Suppose a goal with cost $C$ is validly popped.  If a cheaper solution
of cost $C'<C$ existed, some frontier state on that path would have a current
key at most $C'$.  If unrefined, it would be popped, refined to another lower
bound no greater than $C'$, and requeued; if fully refined, it would precede
the goal.  Either case contradicts the goal being minimal in OPEN.  Finite
maps, positive costs, stale-entry rejection, and reopening give completeness
for every valid query.

Moreover, \EagerFull{}, \LazyFull{}, and \Staged{} all expand non-goal states
only with the same \HFull{} value.  Promotion preserves ordinal, and
Equation~\ref{eq:open} fixes final-stage ordering.  Therefore they should
enumerate successors in the same sequence.  The implementation hashes that
sequence.  Digest equality is used as a fail-closed empirical invariant in
addition to the proof and independent BFS cost check.

Counters distinguish deterministic work from timing: expanded and generated
states, relaxations, pops, stale pops, requeues, calls to Manhattan/prefix/
suffix/full stages, pivot evaluations, table reads, cache hits, discovered
states, and maximum OPEN/live-state sizes.  Timing covers the complete search,
while per-stage timers are disabled to avoid instrumentation changing the
quantity measured.
